\section*{Chapter \ref{ch:m3:refined-products-and-cracks} --- Refined Products and Cracks}

\begin{solution}{exo:m3:refined-products-and-cracks:1}
Gasoline \qty{100.80}{\$/\barrel}, diesel \qty{117.60}{\$/\barrel}; $(2 \times 100.80 +
117.60 - 3 \times 75)/3 = \qty{31.40}{\$/\barrel}$.
\end{solution}

\begin{solution}{exo:m3:refined-products-and-cracks:2}
7.44 barrels a tonne: $900/7.4435 = \qty{120.91}{\$/\barrel}$.
\end{solution}

\begin{solution}{exo:m3:refined-products-and-cracks:3}
Summer driving raises demand, and summer-grade gasoline (RVP at most 9.0 psi from 1
June) needs fewer cheap volatile components and costs more to make. The futures contract
switches to the summer grade in the spring, so the price steps up with the grade.
\end{solution}

\begin{solution}{exo:m3:refined-products-and-cracks:4}
$0.47 \times 100 + 0.30 \times 110 + 0.30 \times 60 - 80 = \qty{18}{\$/\barrel}$. The
yields are volumes: the products are less dense than the crude, so their volume exceeds
the crude's (processing gain).
\end{solution}

\begin{solution}{exo:m3:refined-products-and-cracks:5}
$2.62 - 2.50 - 0.08 = +0.04$ \$/gal: open. It closes when the Gulf price reaches
\qty{2.54}{\$/gal}.
\end{solution}

\begin{solution}{exo:m3:refined-products-and-cracks:6}
1{,}800{,}000 barrels: buy 1{,}800 CL, sell 1{,}200 RB and 600 HO.
\end{solution}

\begin{solution}{exo:m3:refined-products-and-cracks:7}
May 2022, \qty{62.68}{\$/\barrel}; the average over July 2006 to August 2026 is
\qty{21.20}{\$/\barrel}.
\end{solution}

\begin{solution}{exo:m3:refined-products-and-cracks:8}
A refiner is long products and short crude. Buying crude futures adds to its exposure to
the crude price (it will buy crude anyway, at whatever price) instead of hedging it, and
leaves its product sales unhedged. It locks its margin only by buying crude and
\emph{selling} products, in its slate's ratio.
\end{solution}

\begin{solution}{pb:m3:refined-products-and-cracks:1}
\textbf{1.} Gasoline $3.2132 \times 42 = \qty{134.95}{\$/\barrel}$, diesel
\qty{178.63}{\$/\barrel}. \textbf{2.} \qty{65.61}{\$/\barrel}. \textbf{3.} About 3.1 times
the average of \qty{21.20}{}. \textbf{4.} The Gulf disruption of 2026 took crude and
products out of the market, and diesel supply was already tight after Europe lost
Russian imports. \textbf{5.} May 2022, \qty{62.68}{\$/\barrel}.
\textbf{6.} $90\,000 \times 92 = 8.28$ million barrels. \textbf{7.} 8{,}280. \textbf{8.}
5{,}520 RB and 2{,}760 HO. \textbf{9.} $65.61 \times 8.28$ million = USD~543.3 million.
\textbf{10.} The futures gain $(65.61 - 25) \times 8.28$ million = USD~336.3 million,
offsetting the lower margin on the physical barrels.
\textbf{11.} Basis between its crude and WTI Cushing (grade and location). \textbf{12.}
Basis between Gulf and New York products (\cref{fig:m3:refined-products-and-cracks:arb}).
\textbf{13.} Sell products in its own ratio: 4{,}968 RB and 3{,}312 HO. \textbf{14.} The
hedge stays on for barrels the refinery no longer processes: it becomes an open position
long crude and short products, to be closed or rolled. \textbf{15.} Variation margin on
the short product futures when products rise faster than crude, with no cash yet from
the physical margin. \textbf{16.} To keep upside if the crack widens further, to limit
margin calls, and because output is uncertain. \textbf{17.} Speculators, trading houses
and funds long cracks, and refiners' opposite numbers (crude producers selling, product
consumers buying). \textbf{18.} Locking three times the average margin is attractive, but
the curve is usually backwardated at such times and the margin calls on the hedge can be
large. \textbf{19.} \emph{Named result:} buy 8{,}280 CL, sell 5{,}520 RB and 2{,}760 HO,
locking \qty{65.61}{\$/\barrel} of crude, USD~543.3 million for the quarter. \textbf{20.}
The spread between products and crude: the crack, not oil.
\end{solution}

\begin{solution}{iq:m3:refined-products-and-cracks:1}
Three barrels of crude for two of gasoline and one of distillate: roughly the slate of a
US refinery, and all three are liquid futures. The crack is the margin per barrel of
crude.
\iqlookfor{the stylised slate and the futures legs.}
\end{solution}

\begin{solution}{iq:m3:refined-products-and-cracks:2}
Convert with density: barrels per tonne is $6.2898/\rho$ with $\rho$ in kg/l, about 7.45
for gasoil; divide the price per tonne by it.
\iqlookfor{density, not a single magic number.}
\end{solution}

\begin{solution}{iq:m3:refined-products-and-cracks:3}
Regress the crack on calendar-month dummies, or compare month means with a block
bootstrap over years; beware the contract's grade change and the roll. Trade it as a
calendar spread of cracks (long summer against winter months), sized for the years it
fails.
\iqlookfor{a proper test on overlapping years, and a spread rather than an outright.}
\end{solution}

\begin{solution}{iq:m3:refined-products-and-cracks:4}
A distillate-specific supply shock: refinery outages, loss of imports (sanctions,
disrupted routes), a cold winter, a regulation change; gasoline supply is unaffected or
its demand is weak.
\iqlookfor{product-specific supply and demand, not crude.}
\end{solution}

\begin{solution}{iq:m3:refined-products-and-cracks:5}
Crude-quality and location basis, product location and grade basis, slate changes,
volume risk (outages), margin-call liquidity, and the roll of the hedge across months.
\iqlookfor{basis, volume and liquidity.}
\end{solution}

\begin{solution}{iq:m3:refined-products-and-cracks:6}
Reference data per refinery (crude grades and their pricing formulas, slate, costs); a
price store with unit and currency conversions; a margin engine per refinery computing
product value less crude cost; checks on missing prices and stale slates; output by
region and a history for trends.
\iqlookfor{units, per-asset configuration, data quality.}
\end{solution}
